\section{The \ours task} \label{sec:task_definition}
%We formalize our definition of the \ours task and define our evaluation metrics.
\vspace{.2cm} \noindent {\bf Task Formulation} The inputs to our task are a scientific claim $c$ and a corpus of abstracts $\cA$.
All abstracts $a \in \cA$ are labeled as $y(c, a)\in\{$\supports, \refutes, \notenough$\}$ with respect to a claim $c$.  The abstracts that either \support or \refute $c$ are referred to as \emph{evidence abstracts} for $c$, denoted as $\cE(c)$. %, which we denote as $\cE(c) = \{a^1, \dots, a^k\}$.
Each evidence abstract $a \in \cE(c)$ is annotated with rationales. A single rationale $R_i$ is a collection of sentences $\{r_1(c, a), \dots, r_{m}(c, a)\}$, where $m$ is the number of sentences in rationale $R_i$.
We denote the set of all rationales as $\cR(c, a) = \{R_1(c, a), \dots, R_n(c, a)\}$.

Given a claim $c$ and a corpus $\cA$, the system must predict a set of evidence abstracts $\Ehat(c)$. For each abstract $a \in \Ehat(c)$, it must predict a label $\yhat(c, a)$, and a collection of \emph{rationale sentences} $\Shat(c, a) = \{\shat_1(c, a), \dots, \shat_{\ell}(c, a)\}$. Note that although the gold annotations may contain multiple separate rationales, to simplify the prediction task we only require the model to predict a single collection of \emph{rationale sentences}; these sentences may encompass multiple gold rationales.


\vspace{.2cm} \noindent {\bf Task Evaluation} \label{sec:evaluation}
We evaluate the task at two levels of granularity. For \emph{abstract-level} evaluation, we assess the model's ability to identify the abstracts in the corpus that support or refute the claim. For \emph{sentence-level} evaluation, we evaluate the model's performance at identifying the sentences sufficient to justify the abstract-level predictions. We conduct evaluations in both the ``Open'' \fever-style \cite{Thorne2018FEVERAL} setting where the evidence abstracts must be retrieved, and the ``Oracle abstract'' \eraser-style \cite{DeYoung2019ERASERAB} setting where the gold evidence abstracts $\cE(c)$ are provided.


\vspace{.2cm} \noindent {\bf Abstract-level evaluation} is inspired by the \fever score. Given a claim $c$, a predicted evidence abstract $a \in \Ehat(c)$ is \emph{correctly labeled} if (1) $a$ is a gold evidence abstract for $c$, and (2) The predicted label is correct: $\yhat(c, a) = y(c, a)$. It is \emph{correctly rationalized} if, in addition, the predicted rationale sentences contain a gold rationale, i.e., there exists some gold rationale $R_i(c, a) \subseteq \Shat(c, a)$.

Like \fever, which limits the maximum number of predicted rationale sentences to five, \ours limits to three predicted rationale sentences. Overall performance is measured by the micro-F1 of the precision and recall over the correctly-labeled and correctly-rationalized evidence abstracts. We refer to these evaluations as \abstractLabelOnly and \abstractLabelRationale, respectively.

\vspace{.2cm} \noindent {\bf Sentence-level evaluation} measures performance in identifying individual rationale sentences. Unlike the abstract-level metrics, this evaluation penalizes the prediction of extra rationale sentences. A predicted rationale sentence $\shat(c, a)$ is correctly \emph{selected} if (1) It is a member of some gold rationale $R_i(c, a)$ and (2) All other sentences from the same gold rationale $R_i(c, a)$ are among the predicted $\Shat(c, a)$. It is correctly \emph{labeled} if, in addition, the abstract $a$ is correctly labeled: $\yhat(c, a) = y(c, a)$. Overall performance is measured by the micro-F1 of the precision and recall of correctly-selected and correctly-labeled rationale sentences, denoted \sentenceSelectionOnly and \sentenceSelectionLabel, respectively. For sentence-level evaluation, we do not limit the number of predicted rationale sentences, since the evaluation penalizes models that over-predict.
